\section*{Bab \ref{ch:g9:elements-universe-earth} --- Unsur-Unsur di Alam Semesta dan di Bumi}

\begin{solution}{exo:g9:elements-universe-earth:1}
Hidrogen (sekitar \qty{74}{\percent}) dan helium (sekitar
\qty{25}{\percent}).
\end{solution}

\begin{solution}{exo:g9:elements-universe-earth:2}
Oksigen pada keduanya.
\end{solution}

\begin{solution}{exo:g9:elements-universe-earth:3}
Di dalam bintang, dan dalam ledakannya.
\end{solution}

\begin{solution}{exo:g9:elements-universe-earth:4}
Sekitar \qty{3.5}{\percent}.
\end{solution}

\begin{solution}{exo:g9:elements-universe-earth:5}
$46.6 + 27.7 = \qty{74.3}{\percent}$: kira-kira tiga perempat.
\end{solution}

\begin{solution}{exo:g9:elements-universe-earth:6}
Keduanya gas yang sangat ringan: keduanya lolos dari Bumi muda ke
antariksa (dan helium tidak membentuk senyawa yang dapat menahannya di
dalam batuan).
\end{solution}

\begin{solution}{exo:g9:elements-universe-earth:7}
$0.23 \times 70 = \qty{16.1}{kg}$.
\end{solution}

\begin{solution}{exo:g9:elements-universe-earth:8}
\omterm{def:g7:atoms-and-molecules:atom}{Atom} oksigen 16 kali lebih berat daripada \omterm{def:g7:atoms-and-molecules:atom}{atom} hidrogen: \omterm{def:g7:atoms-and-molecules:atom}{atom} oksigen
yang lebih sedikit tetap lebih berat.
\end{solution}

\begin{solution}{exo:g9:elements-universe-earth:9}
Sekitar $0.081 \times 1000 = \qty{81}{kg}$ aluminium.
\end{solution}

\begin{solution}{exo:g9:elements-universe-earth:10}
Oksigen dan kalsium (dan, di luar kedelapan unsur utama kerak, hidrogen,
karbon, dan fosfor terdapat di kerak dalam jumlah kecil).
\end{solution}

\begin{solution}{exo:g9:elements-universe-earth:11}
$40 \times 1.67 \times 10^{-27} = \qty{6.68e-26}{kg}$;
$1.06 / 6.68 \times 10^{-26} \approx 1.6 \times 10^{25}$ \omterm{def:g7:atoms-and-molecules:atom}{atom} kalsium.
\end{solution}

\begin{solution}{exo:g9:elements-universe-earth:12}
\omterm{def:g7:atoms-and-molecules:atom}{Atom} tidak diciptakan dan tidak dimusnahkan oleh \omterm{def:g5:irreversible-changes:chemical-change}{perubahan kimia}: \omterm{def:g7:atoms-and-molecules:atom}{atom}
hanya berpindah dari satu senyawa ke senyawa lain (batuan, gas, gula,
tubuh makhluk hidup). \omterm{def:g7:atoms-and-molecules:atom}{Atom-atom} karbon di Bumi dipakai berulang-ulang.
\end{solution}

\begin{solution}{pb:g9:elements-universe-earth:1}
\textbf{1.} $0.61 \times 70 = \qty{42.7}{kg}$.

\textbf{2.} Karbon $0.23 \times 70 = \qty{16.1}{kg}$; hidrogen
$0.10 \times 70 = \qty{7.0}{kg}$.

\textbf{3.} $61 + 23 + 10 = \qty{94}{\percent}$.

\textbf{4.} Air, \ce{H2O}.

\textbf{5.} $16 \times 1.67 \times 10^{-27} = \qty{2.67e-26}{kg}$.

\textbf{6.} Karbon: $12 \times 1.67 \times 10^{-27} = \qty{2.00e-26}{kg}$;
hidrogen: \qty{1.67e-27}{kg}.

\textbf{7.} 16.

\textbf{8.} Sebuah \omterm{def:g9:inside-the-atom:nucleus}{elektron} kira-kira 1836 kali lebih ringan daripada
sebuah \omterm{def:g9:inside-the-atom:proton}{nukleon}: \omterm{def:g9:inside-the-atom:nucleus}{elektron-elektron} itu kurang dari seperseribu massanya.

\textbf{9.} $7.0 / 1.67 \times 10^{-27} \approx 4.2 \times 10^{27}$
\omterm{def:g7:atoms-and-molecules:atom}{atom} hidrogen.

\textbf{10.} $16.1 / 2.00 \times 10^{-26} \approx 8.0 \times 10^{26}$
\omterm{def:g7:atoms-and-molecules:atom}{atom} karbon.

\textbf{11.} $42.7 / 2.67 \times 10^{-26} \approx 1.60 \times 10^{27}$
\omterm{def:g7:atoms-and-molecules:atom}{atom} oksigen.

\textbf{12.} Hitungan kita, $1.60 \times 10^{27}$, sesuai dengan
perkiraan $1.61 \times 10^{27}$: tubuh seberat \qty{70}{kg} mengandung
sekitar $1.6 \times
10^{27}$ \omterm{def:g7:atoms-and-molecules:atom}{atom} oksigen.
\end{solution}
